% !TeX root = ../0_Manuscript.tex

\section{Body Biasing Injection enhanced practice \ddc}
\label{chap2:goodPractices}
With the platform thoroughly introduced,we can now discuss the BBI practice in the state of the art compared to the enhancements we propose.
\subsection{BBI practice in the state of the art \ddc}
% !TeX root = ../0_Manuscript.tex

%\begin{figure}[H]
%    \centering
%    \includegraphics[width=0.7\textwidth, center]{2_goodPractices/figures/model0Probe.pdf}
%    \caption{Simplified model of a \bbi platform in the state-of-the-art. The model represents the major tools involved in a \bbi platform, with the voltage pulse generator, the \bbi metal probe, the transmission line used to interconnect the generator to the probe, the IC, modeled with a 1 k\textOmega\xspace resistor, and the platform grounding impedance, represented with two 150 \textOmega\xspace resistors. This model allows for fast evaluation and prediction of the \bbi platform macro-behavior.}
%    \label{fig:bbiPracticeBadGnd}
%\end{figure}

\begin{figure}[ht]
    \centering
    \begin{subfigure}{0.45\textwidth}
        \includegraphics[width=1.0\textwidth, center]{2_goodPractices/figures/model0Probe.pdf}
        \caption{Model of a state-of-the-art \bbi setup, with its key components.}
        \label{fig:bbiPracticeBadGnd}
    \end{subfigure}
    \hspace{0.045\textwidth}
    \begin{subfigure}{0.45\textwidth}
        \includegraphics[width=1.0\textwidth, center]{2_goodPractices/figures/SIMPLE_BBI_0.pdf}
        \caption{Simulation results of the state-of-the-art \bbi platform model.}
        \label{fig:bbiPracticeBadGndSignals}
    \end{subfigure}
    \caption{State-of-the-art \bbi platform electrical model and its simulation results.}
    \label{fig:bbiBadGndGlobalFig}
\end{figure}

First and foremost, in order to introduce the proposed enhancements, we will explain the limitations of state of the art BBI practices.
Fig. \ref{fig:bbiPracticeBadGnd} presents an example BBI platform as it is performed in the state of the art.
This schematic only represents the main pieces of equipment, roughly composed of:
\begin{itemize}
    \item A voltage pulse generator
    \item A transmission line
    \item A BBI probe
    \item An IC target
    \item A grounding setup
\end{itemize}
In that case, we want to bring forward two limitations of the platform as it stands.
The first one is that most high voltage generator commercially available require a very specific load impedance to be used in order to meet their specifications, thus the \mbox{50 \textOmega} transmission line.
However, the vast majority of integrated circuits substrates do not present at all times and in every location the exact required impedance for a specific generator.
% !TeX root = ../0_Manuscript.tex

\begin{figure}[H]
    \centering
    \includegraphics[width=15cm, center]{2_goodPractices/figures/SIMPLE_BBI_0.pdf}
    \caption{BBI platform example in the state of the art: simulation}
    \label{fig:bbiPracticeBadGndSignals}
\end{figure}
Therefore, there will inevitably be impedance mismatch effects, such as bounces in the transmission line.
These effects are illustrated in Fig. \ref{fig:bbiPracticeBadGndSignals}.
As one can notice, due to the bounces, there is a 25 \% undershoot on the pulse received by the target (green signal), in addition to a 100 V overshoot.
It is directly echoed on the IC ground current measurement (purple signal).
However, one might wonder why this can be a problem for the practice of BBI.
In fact, because an attacker has to achieve precise timing when injecting faults into an IC